\ctitle{面向新型内存体系结构的高效数据管理机制研究}

% 以下个人信息请在形成正式稿时替换。
\xuehao{待填写}
\schoolcode{10487}
\miji{\hei\textbf{公开}}
\csubjectname{待填写}
\cauthorname{待填写}
\csupervisoronename{{待填写}}
\csupervisoronetitle{教\hspace{1em}授}
\defencedate{待填写} \grantdate{}
\chair{}\firstreviewer{}\secondreviewer{}\thirdreviewer{}

\etitle{Efficient Data Management Mechanisms for Emerging Memory Architectures}
\edegree{Doctor of Philosophy}
\esubject{To Be Completed}
\eauthor{To Be Completed}
\esupervisorone{To Be Completed}
\edate{To Be Completed}

\dctab{\begin{tabular}{|P{0.9cm}|P{1.8cm}|P{1.8cm}|P{5.4cm}|}
  \hline
  &{\hei\textbf{姓名}}&{\hei\textbf{职称}}&{\hei\textbf{单位}}\\
  \hline
  主席&待填写&待填写&待填写\\
  \hline
  委员&待填写&待填写&待填写\\
  \hline
\end{tabular}}

\cabstract{
本文面向持久内存、远程内存与计算—内存—存储解耦架构，研究传统数据管理机制与新型内存硬件特征之间的失配问题。论文拟从数据表示、数据布局和管理功能三个层次展开：研究基于地址重映射的持久内存数据寻址机制、面向解耦内存的多版本数据组织机制，以及面向云原生数据库的分离式缓存管理机制。本文件目前为论文大纲稿，摘要将在正文和实验完成后进一步凝练。
}
\ckeywords{新型内存体系结构；持久内存；解耦内存；数据管理；缓存管理}

\eabstract{
This dissertation studies the mismatch between conventional data-management mechanisms and emerging memory architectures, including persistent memory, remote memory, and compute--memory--storage disaggregation. The planned study spans three levels: data representation, data layout, and management-function placement. This document is an outline; the abstract will be revised after the technical chapters and evaluation are complete.
}
\ekeywords{Emerging memory architectures; persistent memory; disaggregated memory; data management; buffer management}
